%%%%%%%%%%%%%%%%
%commands to use drawings in equations
%%%%%%%%%%%%%%%%
\newdimen\BBoxXMin
\newdimen\BBoxYMin
\newdimen\BBoxXMax
\newdimen\BBoxYMax
\colorlet{bgcolor}{gray!20}
\colorlet{defectcolor}{gray!50}

\definecolor{myPink}{HTML}{FDF5FF}
\definecolor{myRed}{HTML}{eb0707}
\definecolor{myOrange}{HTML}{eb7907}
\definecolor{myGrey}{HTML}{D1D5DB}
\colorlet{myPurple}{violet!70!magenta}
%\definecolor{myPurple}{HTML}{51158C}

% parabolic and weyl defect colours
\definecolor{T-region}{HTML}{E88444} % lovely orange
\definecolor{G-region}{HTML}{B3CCE6} % lovely blue
\definecolor{parabolic-defect}{HTML}{77CCAA} % bluegreen
\definecolor{weyl-defect}{HTML}{888811} % green

\tikzset{
    thick,
    %>=stealth,
    font=\small\sffamily,
    % --- ESTILOS ---
    line_defect/.style={draw=defectcolor, line width=4pt}, 
    skein/.style={thick},
    3D_line/.style={thick},
    2D_line/.style={thick},
    1D_line/.style={thick},
    % Estilo para cruces
    crossing_over/.style={thick,preaction={draw=bgcolor, -, line width=4pt}},
    every node/.style={scale=.9}
}



% 1 full Weyl cat
\newcommand{\FullWeylCatScaled}[2]{%
  \mathord{%
    \tikz[
      baseline=-0.5ex,
      scale=#1,
      every node/.style={transform shape}
    ]{%
      \begin{scope}
        \begin{scope}
          \path[clip] (0,0) circle[radius=6mm];
          \fill[T-region] (0,0) circle[radius=6mm];
          \filldraw[fill=G-region,draw=parabolic-defect,very thick]
            (-6mm,3mm) to[out=0,in=0,looseness=1.6]
            coordinate[pos=.5] (left-junction) (-6mm,-3mm);
          \filldraw[fill=G-region,draw=parabolic-defect,very thick]
            (6mm,3mm) to[out=180,in=180,looseness=1.6]
            coordinate[pos=.5] (right-junction) (6mm,-3mm);

          \draw[very thick,draw=weyl-defect]
            (left-junction) --
            node[pos=.5,above,scale=.7] {$ #2\mathcal{T}_{w_0} $}
            (right-junction);

          \fill (left-junction) circle[radius=1.5pt]
                (right-junction) circle[radius=1.5pt];
        \end{scope}

        %\draw[thick]
        %  ( 90-15:6mm-.5mm) -- ( 90-15:6mm+.5mm)
        %  ( 90+15:6mm-.5mm) -- ( 90+15:6mm+.5mm)
        %  (270-15:6mm-.5mm) -- (270-15:6mm+.5mm)
        %  (270+15:6mm-.5mm) -- (270+15:6mm+.5mm);

         \draw[thick] (0,0) circle[radius=6mm];
         \pgfresetboundingbox
      \path[use as bounding box] (-19.7pt,-19.7pt) rectangle (19.7pt,19.7pt);
      \end{scope}
    }%
  }%
}

\newcommand{\FullWeylCat}{%
  \mathchoice
    {\FullWeylCatScaled{1}{\textstyle}}
    {\FullWeylCatScaled{1}{\textstyle}}
    {\FullWeylCatScaled{0.7}{\scriptstyle}}
    {\FullWeylCatScaled{0.55}{\scriptscriptstyle}}%
}

%%%%%%
%%%%%%%
% 2 left Weyl cat
\newcommand{\LeftWeylCatScaled}[2]{%
  \mathord{%
    \tikz[
      baseline=-0.5ex,
      scale=#1,
      every node/.style={transform shape}
    ]{%
      \begin{scope}
        \fill[G-region]
          (120:6mm) coordinate (-NW)
          arc[start angle=120, delta angle=120, radius=6mm]
          to (0,0) coordinate (-center)
          to cycle;

        \fill[T-region]
          (240:6mm) coordinate (-SW)
          arc[start angle=240, delta angle=240, radius=6mm]
          to (0,0)
          to cycle;

        \path (0:6mm) coordinate (-E);

        % defects
        \draw[very thick,parabolic-defect]
          (-center) to coordinate[pos=.5] (-mid-NW) coordinate[pos=1] (-label-NW) (-NW)
          (-center) to coordinate[pos=.5] (-mid-SW) coordinate[pos=1] (-label-SW) (-SW);

        \draw[very thick,weyl-defect]
          (-center) to
          coordinate[pos=.5] (-mid-E)
          coordinate[pos=.8] (-late-E)
          coordinate[black,pos=1] (-label-E)
          (-E);

        \fill (-center) circle[radius=2pt];

        \draw[thick] (-center) circle[radius=6mm];
      \end{scope}
    }%
  }%
}
\newcommand{\LeftWeylCat}{%
  \mathchoice
    {\LeftWeylCatScaled{1}{\textstyle}}
    {\LeftWeylCatScaled{1}{\textstyle}}
    {\LeftWeylCatScaled{0.7}{\scriptstyle}}
    {\LeftWeylCatScaled{0.55}{\scriptscriptstyle}}%
}


% 3 left Weyl alg G x T^2
\newcommand{\LeftWeylAlgGTTScaled}[2]{%
  \mathord{%
    \tikz[
      baseline=-0.5ex,
      scale=#1,
      every node/.style={transform shape}
    ]{%
      \begin{scope}
        \fill[G-region]
          (120:6mm) coordinate (-NW)
          arc[start angle=120, delta angle=120, radius=6mm]
          to (0,0) coordinate (-center)
          to cycle;

        \fill[T-region]
          (240:6mm) coordinate (-SW)
          arc[start angle=240, delta angle=240, radius=6mm]
          to (0,0)
          to cycle;

        \path (0:6mm) coordinate (-E);

        % defects
        \draw[very thick,parabolic-defect]
          (-center) to coordinate[pos=.5] (-mid-NW) coordinate[pos=1] (-label-NW) (-NW)
          (-center) to coordinate[pos=.5] (-mid-SW) coordinate[pos=1] (-label-SW) (-SW);

        \draw[very thick,weyl-defect]
          (-center) to
          coordinate[pos=.5] (-mid-E)
          coordinate[pos=.8] (-late-E)
          coordinate[black,pos=1] (-label-E)
          (-E);

        \fill (-center) circle[radius=2pt];

        \draw[thick]
    ( 90-15:6mm-.5mm) -- ( 90-15:6mm+.5mm)
    ( 90+15:6mm-.5mm) -- ( 90+15:6mm+.5mm)
    (270-15:6mm-.5mm) -- (270-15:6mm+.5mm)
    (270+15:6mm-.5mm) -- (270+15:6mm+.5mm)
    (180-15:6mm-.5mm) -- (180-15:6mm+.5mm)
    (180+15:6mm-.5mm) -- (180+15:6mm+.5mm)
    ;
  \draw[thick]
    (90-15:6mm) arc[start angle=90-15,end angle=-90+15,radius=6mm]
    (90+15:6mm) arc[start angle=90+15,end angle=180-15,radius=6mm]
    (180+15:6mm) arc[start angle=180+15,end angle=270-15,radius=6mm]
    ;
      \end{scope}
    }%
  }%
}
\newcommand{\LeftWeylAlgGTT}{%
  \mathchoice
    {\LeftWeylAlgGTTScaled{1}{\textstyle}}
    {\LeftWeylAlgGTTScaled{1}{\textstyle}}
    {\LeftWeylAlgGTTScaled{0.7}{\scriptstyle}}
    {\LeftWeylAlgGTTScaled{0.55}{\scriptscriptstyle}}%
}
% 4 left Weyl alg T^2
\newcommand{\LeftWeylAlgTTScaled}[2]{%
  \mathord{%
    \tikz[
      baseline=-0.5ex,
      scale=#1,
      every node/.style={transform shape}
    ]{%
      \begin{scope}
        \fill[G-region]
          (120:6mm) coordinate (-NW)
          arc[start angle=120, delta angle=120, radius=6mm]
          to (0,0) coordinate (-center)
          to cycle;

        \fill[T-region]
          (240:6mm) coordinate (-SW)
          arc[start angle=240, delta angle=240, radius=6mm]
          to (0,0)
          to cycle;

        \path (0:6mm) coordinate (-E);

        % defects
        \draw[very thick,parabolic-defect]
          (-center) to coordinate[pos=.5] (-mid-NW) coordinate[pos=1] (-label-NW) (-NW)
          (-center) to coordinate[pos=.5] (-mid-SW) coordinate[pos=1] (-label-SW) (-SW);

        \draw[very thick,weyl-defect]
          (-center) to
          coordinate[pos=.5] (-mid-E)
          coordinate[pos=.8] (-late-E)
          coordinate[black,pos=1] (-label-E)
          (-E);

        \fill (-center) circle[radius=2pt];

        \draw[thick]
    ( 90-15:6mm-.5mm) -- ( 90-15:6mm+.5mm)
    ( 90+15:6mm-.5mm) -- ( 90+15:6mm+.5mm)
    (270-15:6mm-.5mm) -- (270-15:6mm+.5mm)
    (270+15:6mm-.5mm) -- (270+15:6mm+.5mm)
    ;
  \draw[thick]
    (90-15:6mm) arc[start angle=90-15,end angle=-90+15,radius=6mm]
    (90+15:6mm) arc[start angle=90+15,end angle=270-15,radius=6mm]
    ;
      \end{scope}
    }%
  }%
}
\newcommand{\LeftWeylAlgTT}{%
  \mathchoice
    {\LeftWeylAlgTTScaled{1}{\textstyle}}
    {\LeftWeylAlgTTScaled{1}{\textstyle}}
    {\LeftWeylAlgTTScaled{0.7}{\scriptstyle}}
    {\LeftWeylAlgTTScaled{0.55}{\scriptscriptstyle}}%
}
% 5 right Weyl cat
\newcommand{\RightWeylCatScaled}[2]{%
  \mathord{%
    \tikz[
      baseline=-0.5ex,
      scale=#1,
      every node/.style={transform shape}
    ]{%
      \begin{scope}[xscale=-1]
        \fill[G-region]
          (120:6mm) coordinate (-NW)
          arc[start angle=120, delta angle=120, radius=6mm]
          to (0,0) coordinate (-center)
          to cycle;

        \fill[T-region]
          (240:6mm) coordinate (-SW)
          arc[start angle=240, delta angle=240, radius=6mm]
          to (0,0)
          to cycle;

        \path (0:6mm) coordinate (-E);

        % defects
        \draw[very thick,parabolic-defect]
          (-center) to coordinate[pos=.5] (-mid-NW) coordinate[pos=1] (-label-NW) (-NW)
          (-center) to coordinate[pos=.5] (-mid-SW) coordinate[pos=1] (-label-SW) (-SW);

        \draw[very thick,weyl-defect]
          (-center) to
          coordinate[pos=.5] (-mid-E)
          coordinate[pos=.8] (-late-E)
          coordinate[black,pos=1] (-label-E)
          (-E);

        \fill (-center) circle[radius=2pt];

        \draw[thick] (-center) circle[radius=6mm];
      \end{scope}
    }%
  }%
}
\newcommand{\RightWeylCat}{%
  \mathchoice
    {\RightWeylCatScaled{1}{\textstyle}}
    {\RightWeylCatScaled{1}{\textstyle}}
    {\RightWeylCatScaled{0.7}{\scriptstyle}}
    {\RightWeylCatScaled{0.55}{\scriptscriptstyle}}%
}
% 6 rigt Weyl alg G x T^2
\newcommand{\RightWeylAlgGTTScaled}[2]{%
  \mathord{%
    \tikz[
      baseline=-0.5ex,
      scale=#1,
      every node/.style={transform shape}
    ]{%
      \begin{scope}[xscale=-1]
        \fill[G-region]
          (120:6mm) coordinate (-NW)
          arc[start angle=120, delta angle=120, radius=6mm]
          to (0,0) coordinate (-center)
          to cycle;

        \fill[T-region]
          (240:6mm) coordinate (-SW)
          arc[start angle=240, delta angle=240, radius=6mm]
          to (0,0)
          to cycle;

        \path (0:6mm) coordinate (-E);

        % defects
        \draw[very thick,parabolic-defect]
          (-center) to coordinate[pos=.5] (-mid-NW) coordinate[pos=1] (-label-NW) (-NW)
          (-center) to coordinate[pos=.5] (-mid-SW) coordinate[pos=1] (-label-SW) (-SW);

        \draw[very thick,weyl-defect]
          (-center) to
          coordinate[pos=.5] (-mid-E)
          coordinate[pos=.8] (-late-E)
          coordinate[black,pos=1] (-label-E)
          (-E);

        \fill (-center) circle[radius=2pt];

        \draw[thick]
    ( 90-15:6mm-.5mm) -- ( 90-15:6mm+.5mm)
    ( 90+15:6mm-.5mm) -- ( 90+15:6mm+.5mm)
    (270-15:6mm-.5mm) -- (270-15:6mm+.5mm)
    (270+15:6mm-.5mm) -- (270+15:6mm+.5mm)
    (180-15:6mm-.5mm) -- (180-15:6mm+.5mm)
    (180+15:6mm-.5mm) -- (180+15:6mm+.5mm)
    ;
  \draw[thick]
    (90-15:6mm) arc[start angle=90-15,end angle=-90+15,radius=6mm]
    (90+15:6mm) arc[start angle=90+15,end angle=180-15,radius=6mm]
    (180+15:6mm) arc[start angle=180+15,end angle=270-15,radius=6mm]
    ;
      \end{scope}
    }%
  }%
}
\newcommand{\RightWeylAlgGTT}{%
  \mathchoice
    {\RightWeylAlgGTTScaled{1}{\textstyle}}
    {\RightWeylAlgGTTScaled{1}{\textstyle}}
    {\RightWeylAlgGTTScaled{0.7}{\scriptstyle}}
    {\RightWeylAlgGTTScaled{0.55}{\scriptscriptstyle}}%
}

% 7 right Weyl alg T^2
\newcommand{\RightWeylAlgTTScaled}[2]{%
  \mathord{%
    \tikz[
      baseline=-0.5ex,
      scale=#1,
      every node/.style={transform shape}
    ]{%
      \begin{scope}[xscale=-1]
        \fill[G-region]
          (120:6mm) coordinate (-NW)
          arc[start angle=120, delta angle=120, radius=6mm]
          to (0,0) coordinate (-center)
          to cycle;

        \fill[T-region]
          (240:6mm) coordinate (-SW)
          arc[start angle=240, delta angle=240, radius=6mm]
          to (0,0)
          to cycle;

        \path (0:6mm) coordinate (-E);

        % defects
        \draw[very thick,parabolic-defect]
          (-center) to coordinate[pos=.5] (-mid-NW) coordinate[pos=1] (-label-NW) (-NW)
          (-center) to coordinate[pos=.5] (-mid-SW) coordinate[pos=1] (-label-SW) (-SW);

        \draw[very thick,weyl-defect]
          (-center) to
          coordinate[pos=.5] (-mid-E)
          coordinate[pos=.8] (-late-E)
          coordinate[black,pos=1] (-label-E)
          (-E);

        \fill (-center) circle[radius=2pt];

        \draw[thick]
    ( 90-15:6mm-.5mm) -- ( 90-15:6mm+.5mm)
    ( 90+15:6mm-.5mm) -- ( 90+15:6mm+.5mm)
    (270-15:6mm-.5mm) -- (270-15:6mm+.5mm)
    (270+15:6mm-.5mm) -- (270+15:6mm+.5mm)
    ;
  \draw[thick]
    (90-15:6mm) arc[start angle=90-15,end angle=-90+15,radius=6mm]
    (90+15:6mm) arc[start angle=90+15,end angle=270-15,radius=6mm]
    ;
      \end{scope}
    }%
  }%
}
\newcommand{\RightWeylAlgTT}{%
  \mathchoice
    {\RightWeylAlgTTScaled{1}{\textstyle}}
    {\RightWeylAlgTTScaled{1}{\textstyle}}
    {\RightWeylAlgTTScaled{0.7}{\scriptstyle}}
    {\RightWeylAlgTTScaled{0.55}{\scriptscriptstyle}}%
}
% 8 line weyl cat
\newcommand{\WeylCatScaled}[2]{%
  \mathord{%
    \tikz[
      baseline=-0.5ex,
      scale=#1,
      every node/.style={transform shape}
    ]{%
\begin{scope}% clip
    \path[clip] (0,0) circle[radius=6mm];
    \fill[T-region] (0,0) circle[radius=6mm];
  \draw[very thick,draw=weyl-defect]
    (-6mm,0) -- (6mm,0);
\end{scope}

  % boundary
  \draw[thick] (0,0) circle[radius=6mm];
  \pgfresetboundingbox
      \path[use as bounding box] (-19.7pt,-19.7pt) rectangle (19.7pt,19.7pt);
  
    }%
  }%
}
\newcommand{\WeylCat}{%
  \mathchoice
    {\WeylCatScaled{1}{\textstyle}}
    {\WeylCatScaled{1}{\textstyle}}
    {\WeylCatScaled{0.7}{\scriptstyle}}
    {\WeylCatScaled{0.55}{\scriptscriptstyle}}%
}
% 9 line weyl alg T^2
\newcommand{\WeylAlgScaled}[2]{%
  \mathord{%
    \tikz[
      baseline=-0.5ex,
      scale=#1,
      every node/.style={transform shape}
    ]{%
\begin{scope}% clip
    \path[clip] (0,0) circle[radius=6mm];
    \fill[T-region] (0,0) circle[radius=6mm];
    
  \end{scope}
  \draw[very thick,draw=weyl-defect]
    (-6mm,0) -- (6mm,0);


  % boundary
  \draw[thick]
    ( 90-15:6mm-.5mm) -- ( 90-15:6mm+.5mm)
    ( 90+15:6mm-.5mm) -- ( 90+15:6mm+.5mm)
    (270-15:6mm-.5mm) -- (270-15:6mm+.5mm)
    (270+15:6mm-.5mm) -- (270+15:6mm+.5mm)
    ;
  \draw[thick]
    (90-15:6mm) arc[start angle=90-15,end angle=-90+15,radius=6mm]
    (90+15:6mm) arc[start angle=90+15,end angle=270-15,radius=6mm]
    ;
    }%
  }%
}
\newcommand{\WeylAlg}{%
  \mathchoice
    {\WeylAlgScaled{1}{\textstyle}}
    {\WeylAlgScaled{1}{\textstyle}}
    {\WeylAlgScaled{0.7}{\scriptstyle}}
    {\WeylAlgScaled{0.55}{\scriptscriptstyle}}%
}
% 10 full weyl alg T^2
\newcommand{\FullWeylAlgScaled}[2]{%
  \mathord{%
    \tikz[
      baseline=-0.5ex,
      scale=#1,
      every node/.style={transform shape}
    ]{%
      \begin{scope}
        \begin{scope}
          \path[clip] (0,0) circle[radius=6mm];
          \fill[T-region] (0,0) circle[radius=6mm];
          \filldraw[fill=G-region,draw=parabolic-defect,very thick]
            (-6mm,3mm) to[out=0,in=0,looseness=1.6]
            coordinate[pos=.5] (left-junction) (-6mm,-3mm);
          \filldraw[fill=G-region,draw=parabolic-defect,very thick]
            (6mm,3mm) to[out=180,in=180,looseness=1.6]
            coordinate[pos=.5] (right-junction) (6mm,-3mm);

          \draw[very thick,draw=weyl-defect]
            (left-junction) --
            (right-junction);

          \fill (left-junction) circle[radius=1.5pt]
                (right-junction) circle[radius=1.5pt];
        \end{scope}

        \draw[thick]
          ( 90-15:6mm-.5mm) -- ( 90-15:6mm+.5mm)
          ( 90+15:6mm-.5mm) -- ( 90+15:6mm+.5mm)
          (270-15:6mm-.5mm) -- (270-15:6mm+.5mm)
          (270+15:6mm-.5mm) -- (270+15:6mm+.5mm);

        \draw[thick]
          (90-15:6mm) arc[start angle=90-15,end angle=-90+15,radius=6mm]
          (90+15:6mm) arc[start angle=90+15,end angle=270-15,radius=6mm];
      \end{scope}
    }%
  }%
}

\newcommand{\FullWeylAlg}{%
  \mathchoice
    {\FullWeylAlgScaled{1}{\textstyle}}
    {\FullWeylAlgScaled{1}{\textstyle}}
    {\FullWeylAlgScaled{0.7}{\scriptstyle}}
    {\FullWeylAlgScaled{0.55}{\scriptscriptstyle}}%
}

%Digon

\newcommand{\DigonScaled}[2]{%
  \mathord{%
    \tikz[
      baseline=-0.5ex,
      scale=#1,
      every node/.style={transform shape}
    ]{%
       \begin{scope} % clip
    \path[clip] (0,0) circle[radius=6mm];
    \fill[G-region] (0,0) circle[radius=6mm];
    \fill[T-region] (-6mm,3mm) -- ++(16mm,0) -- ++ (0,3mm) -- ++ (-16mm,0);
    \fill[T-region] (-6mm,-3mm) -- ++(16mm,0) -- ++ (0,-3mm) -- ++ (-20mm,0);

    \draw[very thick,parabolic-defect] (-6mm,3mm) --node[black,pos=.3,below,scale=.7] {$\widetilde{\mathcal{D}_B}$} (6mm,3mm) (-6mm,-3mm) -- node[black,pos=.7,above,scale=.7] {$\widetilde{\mathcal{D}_B}$} (6mm,-3mm);
  \end{scope}
   \draw[thick] (0,0) circle[radius=6mm];
   \pgfresetboundingbox
      \path[use as bounding box] (-19.7pt,-19.7pt) rectangle (19.7pt,19.7pt);
    }%
  }%
}

\newcommand{\Digon}{%
  \mathchoice
    {\DigonScaled{1}{\textstyle}}
    {\DigonScaled{1}{\textstyle}}
    {\DigonScaled{0.7}{\scriptstyle}}
    {\DigonScaled{0.55}{\scriptscriptstyle}}%
}

%DigonNoLabel
%Digon

\newcommand{\DigonNoLabelScaled}[2]{%
  \mathord{%
    \tikz[
      baseline=-0.5ex,
      scale=#1,
      every node/.style={transform shape}
    ]{%
       \begin{scope} % clip
    \path[clip] (0,0) circle[radius=6mm];
    \fill[G-region] (0,0) circle[radius=6mm];
    \fill[T-region] (-6mm,3mm) -- ++(16mm,0) -- ++ (0,3mm) -- ++ (-16mm,0);
    \fill[T-region] (-6mm,-3mm) -- ++(16mm,0) -- ++ (0,-3mm) -- ++ (-20mm,0);

   \draw[very thick,parabolic-defect] (-6mm,3mm) --(6mm,3mm) (-6mm,-3mm) --  (6mm,-3mm);
  \end{scope}
   \draw[thick] (0,0) circle[radius=6mm];
   \pgfresetboundingbox
      \path[use as bounding box] (-19.7pt,-19.7pt) rectangle (19.7pt,19.7pt);
    }%
  }%
}

\newcommand{\DigonNoLabel}{%
  \mathchoice
    {\DigonNoLabelScaled{1}{\textstyle}}
    {\DigonNoLabelScaled{1}{\textstyle}}
    {\DigonNoLabelScaled{0.7}{\scriptstyle}}
    {\DigonNoLabelScaled{0.55}{\scriptscriptstyle}}%
}
%DigonAlg
\newcommand{\DigonAlgScaled}[2]{%
  \mathord{%
    \tikz[
      baseline=-0.5ex,
      scale=#1,
      every node/.style={transform shape}
    ]{%
       \begin{scope} % clip
    \path[clip] (0,0) circle[radius=6mm];
    \fill[G-region] (0,0) circle[radius=6mm];
    \fill[T-region] (-6mm,3mm) -- ++(16mm,0) -- ++ (0,3mm) -- ++ (-16mm,0);
    \fill[T-region] (-6mm,-3mm) -- ++(16mm,0) -- ++ (0,-3mm) -- ++ (-20mm,0);

    \draw[very thick,parabolic-defect] (-6mm,3mm) --node[black,pos=.3,below,scale=.7] {} (6mm,3mm) (-6mm,-3mm) -- node[black,pos=.7,above,scale=.7] {} (6mm,-3mm);
  \end{scope}
   \draw[thick]
          ( 90-15:6mm-.5mm) -- ( 90-15:6mm+.5mm)
          ( 90+15:6mm-.5mm) -- ( 90+15:6mm+.5mm)
          (270-15:6mm-.5mm) -- (270-15:6mm+.5mm)
          (270+15:6mm-.5mm) -- (270+15:6mm+.5mm);

        \draw[thick]
          (90-15:6mm) arc[start angle=90-15,end angle=-90+15,radius=6mm]
          (90+15:6mm) arc[start angle=90+15,end angle=270-15,radius=6mm];
   \pgfresetboundingbox
      \path[use as bounding box] (-19.7pt,-19.7pt) rectangle (19.7pt,19.7pt);
    }%
  }%
}

\newcommand{\DigonAlg}{%
  \mathchoice
    {\DigonAlgScaled{1}{\textstyle}}
    {\DigonAlgScaled{1}{\textstyle}}
    {\DigonAlgScaled{0.7}{\scriptstyle}}
    {\DigonAlgScaled{0.55}{\scriptscriptstyle}}%
}

%G-bubble
\newcommand{\GbubbleScaled}[1]{%
  \mathord{%
    \tikz[
      baseline=-0.5ex,
      scale=#1,
      every node/.style={transform shape}
    ]{%
      % Background
      \fill[T-region] (0,0) circle[radius=6mm];

      % Junction coordinates
      \coordinate (left-junction)  at (180:2mm);
      \coordinate (right-junction) at (0:2mm);

      % Lines from the boundary to the bubble
      \draw[very thick,parabolic-defect]
        (180:6mm) -- (left-junction)
        (0:6mm)   -- (right-junction);

      % Inner bubble
      \filldraw[
        very thick,
        draw=parabolic-defect,
        fill=G-region
      ]
        (right-junction)
          to[out=90,in=90,looseness=1.6]
        (left-junction)
          to[out=-90,in=-90,looseness=1.6]
        cycle;

      % Junction dots
      \fill
        (right-junction) circle[radius=1.5pt]
        (left-junction)  circle[radius=1.5pt];

      % Outer boundary
      \draw[thick] (0,0) circle[radius=6mm];

      % Consistent bounding box
      \pgfresetboundingbox
      \path[use as bounding box]
        (-19.7pt,-19.7pt) rectangle (19.7pt,19.7pt);
    }%
  }%
}

\newcommand{\Gbubble}{%
  \mathchoice
    {\GbubbleScaled{1}}
    {\GbubbleScaled{1}}
    {\GbubbleScaled{0.7}}
    {\GbubbleScaled{0.55}}%
}

%GbubbleAlg
%G-bubble
\newcommand{\GbubbleAlgScaled}[1]{%
  \mathord{%
    \tikz[
      baseline=-0.5ex,
      scale=#1,
      every node/.style={transform shape}
    ]{%
      % Background
      \fill[T-region] (0,0) circle[radius=6mm];

      % Junction coordinates
      \coordinate (left-junction)  at (180:2mm);
      \coordinate (right-junction) at (0:2mm);

      % Lines from the boundary to the bubble
      \draw[very thick,parabolic-defect]
        (180:6mm) -- (left-junction)
        (0:6mm)   -- (right-junction);

      % Inner bubble
      \filldraw[
        very thick,
        draw=parabolic-defect,
        fill=G-region
      ]
        (right-junction)
          to[out=90,in=90,looseness=1.6]
        (left-junction)
          to[out=-90,in=-90,looseness=1.6]
        cycle;

      % Junction dots
      \fill
        (right-junction) circle[radius=1.5pt]
        (left-junction)  circle[radius=1.5pt];

      % Outer boundary
      \draw[thick]
          ( 90-15:6mm-.5mm) -- ( 90-15:6mm+.5mm)
          ( 90+15:6mm-.5mm) -- ( 90+15:6mm+.5mm)
          (270-15:6mm-.5mm) -- (270-15:6mm+.5mm)
          (270+15:6mm-.5mm) -- (270+15:6mm+.5mm);

        \draw[thick]
          (90-15:6mm) arc[start angle=90-15,end angle=-90+15,radius=6mm]
          (90+15:6mm) arc[start angle=90+15,end angle=270-15,radius=6mm];

      % Consistent bounding box
      \pgfresetboundingbox
      \path[use as bounding box]
        (-19.7pt,-19.7pt) rectangle (19.7pt,19.7pt);
    }%
  }%
}

\newcommand{\GbubbleAlg}{%
  \mathchoice
    {\GbubbleAlgScaled{1}}
    {\GbubbleAlgScaled{1}}
    {\GbubbleAlgScaled{0.7}}
    {\GbubbleAlgScaled{0.55}}%
}







 



